% !PRINTED 511
\weberpage{511}

Wir nennen daher einen Theiler $Q$ der Gruppe $P$, der diese
Eigenschaft hat, einen \gesperrt{Normaltheiler} (oder normalen Theiler)\footnote{\gesperrt{Galois} spricht von der eigentlichen Zerlegung (d\'ecomposition propre) einer Gruppe; daher der Ausdruck \glqq{}eigentliche Theiler\grqq{} der auch im Gebrauch ist; neuere Schriftsteller bezeichnen die Normaltheiler als \glqq{}ausgezeichnete oder invariante Untergruppen\grqq{}.}.

Ist aber $Q$ ein Normaltheiler, so ist die Adjunction einer
einzigen Wurzel $\psi$ der H\"ulfsgleichung gleichbedeutend mit der
gleichzeitigen Adjunction aller dieser Wurzeln.

5. \gesperrt{Ist $Q$ selbst nicht normal, so ist der gr\"osste
gemeinschaftliche Theiler $R$ aller conjugirten
Theiler $\pi^{-1}Q\pi$ gewiss normal.}

Denn ist eine Permutation $\varkappa$ in $R$, also in allen Gruppen
$\pi^{-1}Q\pi$ enthalten, so gilt das Gleiche auch von jedem $\pi^{-1}\varkappa\pi$,
d. h. $\pi^{-1}R\pi$ ist Theiler von $R$, und folglich, da es von gleichem
Grade ist, $= R$.

Eine Gruppe, die ausser sich selbst und der identischen
Gruppe keinen Normaltheiler hat, heisst eine \gesperrt{einfache Gruppe}.
Der Durchschnitt aller mit einer einfachen Gruppe conjugirten
Gruppen ist entweder die Gruppe selbst oder die identische Gruppe.

Ist $Q$ ein Normaltheiler von $P$, und $R$ ein Normaltheiler
von $Q$, so ist zwar $R$ ein Theiler von $P$, aber keineswegs immer
normal. Dagegen ist, wenn $R$ ein Normaltheiler von $P$ ist, $R$
auch Normaltheiler von jedem Theiler $Q$ von $P$, von dem $R$
Theiler ist.

\begin{center}
\S.\,156.

\gesperrt{Die Gruppe der Resolventen.}
\end{center}

Die H\"ulfsgleichung $\varphi(t) = 0$, von der die Bestimmung der
Function $\psi$ abh\"angt, geht in die \gesperrt{Galois}'sche Resolvente \"uber,
wenn der Theiler $Q$, zu dem $\psi$ geh\"ort, die identische Gruppe
ist. Denn dann kann nach dem Satze von \gesperrt{Lagrange} (\S.\,155, 2.)
jede Function des K\"orpers $N$, also auch die Wurzeln $\alpha$ selbst
rational durch $\psi$ ausgedr\"uckt werden, und $N$ ist mit $\Omega(\psi)$
identisch.

Wir wollen diese Gleichungen $\varphi(t) = 0$ daher in einem allgemeinen Sinne \gesperrt{Resolventen} nennen. Es sind aber hier zwei
F\"alle zu unterscheiden.
